\item Benjamín y Jacobo echan un vistazo a un acuario que tiene un frente curvo de plástico con espesor uniforme y un radio de curvatura de magnitud $R = 2.25\text{ m}$.
\begin{enumerate}
    \item Localice las imágenes de los peces que se encuentran a (i) $5.00\text{ cm}$ y (ii) $25.0\text{ cm}$ de la pared frontal del acuario.
    \item Encuentre la magnificación de las imágenes (i) y (ii) del inciso (a).
    \item Explique por qué no necesita saber el índice de refracción del plástico para resolver este problema.
    \item Si este acuario es muy largo desde el frente hacia atrás, ¿la imagen de un pez puede estar cada vez más lejos de la superficie frontal que los peces en sí?
    \item En caso negativo, explique por qué no. Si es así, dé un ejemplo y encuentre la magnificación correspondiente.
\end{enumerate}